\clearpage
\begin{summarybox}
\textbf{\textcolor{CSummary}{快记：一个周期 + 一个对称位置}}
\[
\boxed{\begin{gathered}
\text{轴：}\quad x=a+\frac{kT}{2}\\[2pt]
\text{中心：}\quad\left(b+\frac{kT}{2},\,c\right)
\end{gathered}\qquad(k\in\mathbb Z)}
\]
两者都只保证\textbf{这一组一定存在}；已知最小正周期 $T_0$ 才能写出全部对称轴 $x=a+kT_0/2$。

\medskip
\textbf{\textcolor{CSummary}{快懂：为什么是半个周期？}}

\textbf{重复}表示“水平移动整段后函数值不变”；\textbf{轴对称}表示“折痕两侧等距点值相同”；\textbf{中心对称}表示“两侧等距点的值相加为中心高度的两倍”。将一次重复与一次对称结合，就能把已知位置的左右关系搬到半周期以外。

\medskip
\textbf{\textcolor{CSummary}{快用：考试时三步走}}
\begin{enumerate}[leftmargin=*,itemsep=0.3em]
\item \textbf{看条件：}有周期？有轴还是有中心？给的是“一个”还是“最小”？
\item \textbf{选公式：}一轴 + 周期用 $a+kT/2$；两轴用 $2|a_1-a_2|$；一轴一中心用 $f(x+2(b-a))=2c-f(x)$。
\item \textbf{核对范围：}能否断言“全部”？能否断言“最小”？中心的纵坐标是否保留？
\end{enumerate}

\medskip
\textbf{\textcolor{CSummary}{两条值得记住的额外规律}}
\[
\boxed{\begin{aligned}
&\text{两条对称轴}\ \Longrightarrow\ 2|a_1-a_2|\text{ 是周期},\\
&\text{轴 $x=a$ + 中心 $(b,c)$}\ \Longrightarrow\
 f(x+2(b-a))=2c-f(x).
\end{aligned}}
\]
如果还有最小正周期 $T_0$，则一轴一中心的横向距离满足
\[
\boxed{b-a=\frac{(2k+1)T_0}{4}.}
\]

\medskip
\textbf{\textcolor{CSummary}{一张脑海里的图}}
\[
\underbrace{\text{对称轴}}_{\text{波峰}}
\ \xrightarrow{\ T_0/4\ }\
\underbrace{\text{对称中心}}_{\text{中间高度}}
\ \xrightarrow{\ T_0/4\ }\
\underbrace{\text{对称轴}}_{\text{波谷}}
\]
上面表示常见的最近邻情形；一般情况下，要记得可能跨过多段。

\medskip
\textbf{\textcolor{CSummary}{最后只背三句话}}
\begin{center}
\textbf{半周期出现同类对称；两倍轴距产生周期；四分之一周期连着轴和中心。}

\smallskip
\textbf{“一个周期”不等于“最小正周期”；“能找到一组”不等于“全部找齐”。}
\end{center}
\end{summarybox}
